
\section*{\SGAInlineRefTarget{sga3:xxiv:section:4}4. Isotriviality of reductive groups and of principal bundles under reductive groups}

\subsection*{\SGAInlineRefTarget{sga3:xxiv:subsection:4-1}4.1. Definitions. Isotriviality theorem}

\textbf{\SGAInlineRefTarget{sga3:xxiv:definition:4-1-1}Definition 4.1.1.} \emph{Let \(S\) be a scheme, \(G\) an \(S\)-group scheme, \(P\) a principal homogeneous bundle under \(G\). One
says that \(P\) is} locally isotrivial \emph{(resp.} semi-locally isotrivial) if for every point \(s \in S\) (resp. every
finite set \(F\) of points of \(S\) contained in an affine open set) there exist an open set \(U\) of \(S\) containing \(s\)
(resp. \(F\)) and a finite surjective étale morphism \(S' \to U\) such that \(P_{S'}\) is trivial.

\textbf{\SGAInlineRefTarget{sga3:xxiv:definition:4-1-2}Definition 4.1.2.} \emph{Let \(S\) be a scheme, \(G\) an \(S\)-reductive group. One says that \(G\) is} locally isotrivial
\emph{(resp.} semi-locally isotrivial) if for every point \(s \in S\) (resp. every finite set \(F\) of points of \(S\) contained
in an affine open set) there exist an open set \(U\) of \(S\) containing \(s\) (resp. \(F\)) and a finite surjective étale
morphism \(S' \to U\) such that \(G_{S'}\) is splittable.*

\textbf{\SGAInlineRefTarget{sga3:xxiv:remark:4-1-3}Remark 4.1.3.} (i) The equivalence of categories \SGARefLink{sga3:xxiv:definition:4-1-1}{4.1.1} of \SGARefLink{sga3:xxiv:1.17}{1.17} respects by definition local (resp. semi-local)
isotriviality.

\begin{enumerate}
\def\labelenumi{(\roman{enumi})}
\setcounter{enumi}{1}
\item
  Add to the conditions of \SGARefLink{sga3:xxiv:definition:4-1-1}{4.1.1}: \(G\) locally of finite presentation over \(S\). Then the principal bundle \(P\) (or the
  reductive group \(G\)) is locally isotrivial (resp. semi-locally isotrivial) if and only if for every \(S' \to S\), \(S'\)
  local (resp. semi-local), \(P_{S'}\) is isotrivial (or \(G_{S'}\) isotrivial), that is to say, if there exists \(S'' \to S'\)
  finite surjective étale such that \(P_{S''}\) is trivial (or \(G_{S''}\) split).
\item
  In the case of tori, \SGARefLink{sga3:xxiv:definition:4-1-2}{definition 4.1.2} coincides with that of \SGARefLink{sga3:ix:definition:1:1}{Exp. IX, 1.1}.
\end{enumerate}

\subsubsection*{4.1.4.}

Recall (\SGARefLink{sga3:xxii:section:4.3}{Exp. XXII, 4.3} and \SGARefLink{sga3:xxiv:lemma:6-2}{6.2}) that for every reductive group \(G\), we have introduced the groups \(rad(G)\), \(corad(G)\)
and \(d\acute{e}r(G)\). The groups \(rad(G)\) and \(corad(G)\) are tori, which are split when \(G\) is split; moreover, there
exists an isogeny \(rad(G) \to corad(G)\). The \(S\)-group \(d\acute{e}r(G)\) is semisimple, one has \(G / d\acute{e}r(G) =
corad(G)\); it follows that for every principal homogeneous bundle \(P\) under \(G\), \(P / d\acute{e}r(G)\) is a principal
homogeneous bundle under \(corad(G)\). This said, one has:

\textbf{\SGAInlineRefTarget{sga3:xxiv:theorem:4-1-5}Theorem 4.1.5.} \emph{Let \(S\) be a scheme, \(G\) an \(S\)-reductive group of constant type.}

\emph{\SGAInlineRefTarget{sga3:xxiv:theorem:4-1-5:item:i}(i) The following conditions are equivalent:}

\emph{\SGAInlineRefTarget{sga3:xxiv:theorem:4-1-5:item:a}(a) \(G\) is locally (resp. semi-locally) isotrivial.}

\emph{\SGAInlineRefTarget{sga3:xxiv:theorem:4-1-5:item:b}(b) The torus \(rad(G)\) is.}

\emph{(b') The torus \(corad(G)\) is.}

\emph{\SGAInlineRefTarget{sga3:xxiv:theorem:4-1-5:item:c}(c) The Galois covering of \(S\) associated to \(G\) (\SGARefLink{sga3:xxiv:3.11}{3.11}) is.}

\emph{If \(T\) is a maximal torus of \(G\), these conditions are also equivalent to}

\emph{\SGAInlineRefTarget{sga3:xxiv:theorem:4-1-5:item:d}(d) The torus \(T\) is locally (resp. semi-locally) isotrivial.}

\emph{\SGAInlineRefTarget{sga3:xxiv:theorem:4-1-5:item:ii}(ii) Let \(P\) be a principal homogeneous bundle under \(G\); for \(P\) to be locally (resp. semi-locally) isotrivial, it is
necessary and sufficient that the \(corad(G)\)-principal bundle \(P / d\acute{e}r(G)\) be so.}

\textbf{\SGAInlineRefTarget{sga3:xxiv:corollary:4-1-6}Corollary 4.1.6.} \emph{The conditions of \SGARefLink{sga3:xxiv:theorem:4-1-5:item:i}{(i)} are satisfied when \(G\) is semisimple or when \(S\) is locally noetherian and
normal (or more generally geometrically unibranch). The conditions of \SGARefLink{sga3:xxiv:theorem:4-1-5:item:ii}{(ii)} are satisfied when \(G\) is locally (resp.
semi-locally) isotrivial.}

For \SGARefLink{sga3:xxiv:theorem:4-1-5:item:i}{(i)}, the first assertion is trivial from \SGARefLink{sga3:xxiv:theorem:4-1-5:item:b}{(b)}, the second follows from \SGARefLink{sga3:xxiv:theorem:4-1-5:item:c}{(c)} and \SGARefLink{sga3:x:corollary:5-14}{Exp. X, 5.14} and 5.15. For \SGARefLink{sga3:xxiv:theorem:4-1-5:item:ii}{(ii)}, it
suffices to note that by virtue of \SGARefLink{sga3:viii:corollary:4:5}{theorem 90}, a principal bundle under a split torus is semi-locally isotrivial.

\subsection*{\SGAInlineRefTarget{sga3:xxiv:subsection:4-2}4.2. Proof: the semisimple case}

Let us first prove, for later reference:

\textbf{\SGAInlineRefTarget{sga3:xxiv:proposition:4-2-1}Proposition 4.2.1.} \emph{Let \(S\) be a scheme, \(G\) an \(S\)-reductive group, \(T\) a maximal torus of \(G\) (resp. \(B\) a Borel
subgroup, resp. \(B \supset T\) a Killing couple of \(G\)), \(P\) a principal homogeneous bundle under \(G\), \(G'\) the twisted
form of \(G\) associated to \(P\) (via the morphism \(int : G \to \operatorname{Aut}_{S-gr.}(G)\)). One has canonical
isomorphisms}

\[
  P/\operatorname{Norm}_G(T)\xrightarrow{\sim}\operatorname{Tor}(G'),
  \qquad
  P/B\xrightarrow{\sim}\operatorname{Bor}(G'),
  \qquad
  P/T\xrightarrow{\sim}\operatorname{Kil}(G').
\]

By construction, \(G'\) is the quotient of \(P \times_{S} G\) by a certain action of \(G\) (\(((p, g')) g = (pg, g^{-1} g'
g)\)); one therefore has a morphism \(P \times_{S} G \to G'\), that is, a morphism

\[ P \to \operatorname{Hom}_{S}(G, G'), \]

which, as one sees at once, factors through a morphism

\[
  f:P\longrightarrow\operatorname{Isom}_{S\text{-}\mathrm{gr.}}(G,G').
\]

(to verify this assertion, one may assume \(G = P\), in which case one has \(G' = G\), \(f = int\)). The set-theoretic map \(p
\mapsto f(p)(T)\) defines a morphism

\[ P \to Tor(G'). \]

To verify that this morphism induces an isomorphism \(P / Norm_{G}(T) \xrightarrow{\sim} Tor(G')\) as announced, one may
again assume \(P = G\) in which case one is reduced to \SGARefLink{sga3:xxii:cor:5.8.3:item:iii}{Exp. XXII, 5.8.3 (iii)}. (In fact, \emph{loc. cit.} should be replaced by
the statement above.) One argues similarly for \(\operatorname{Bor}\) and \(\operatorname{Kil}\).

\textbf{\SGAInlineRefTarget{sga3:xxiv:proposition:4-2-2}Proposition 4.2.2.} \emph{Let \(S\) be a semi-local scheme, \(G\) an \(S\)-semisimple group of constant type.}

\emph{\SGAInlineRefTarget{sga3:xxiv:proposition:4-2-2:item:i}(i) \(G\) is isotrivial.}

\emph{(ii) Every principal bundle under \(G\) is isotrivial.}

Let us prove \SGARefLink{sga3:xxiv:proposition:4-2-2:item:i}{(i)}. Up to making a finite surjective étale extension of the base, one may, by \SGARefLink{sga3:xxiv:3.9.2}{3.9.2}, assume \(G\)
quasi-split. But then \(G\) is isotrivial by construction (\SGARefLink{sga3:xxiv:3.11}{3.11}, the group \(E\) being finite). To prove \SGARefLink{sga3:xxiv:3.11:item:ii}{(ii)}, one may, by
\SGARefLink{sga3:xxiv:3.11:item:i}{(i)}, assume \(G\) split; one is then reduced to:

\textbf{\SGAInlineRefTarget{sga3:xxiv:lemma:4-2-3}Lemma 4.2.3.} \emph{Let \(S\) be a semi-local scheme. Every principal bundle under a split reductive group is isotrivial.}

Indeed, with the notation of \SGARefLink{sga3:xxiv:proposition:4-2-1}{4.2.1}, where \(B \supset T\) denotes the canonical Killing couple of the split \(G\), the
\(S\)-scheme \(Kil(G')\) possesses a section, after finite surjective étale extension of the base, by \SGARefLink{sga3:xxiv:3.9.2}{3.9.2}. One may
therefore reduce the structure group of \(G\) to \(T\); but \(T\) is split, so \(H^{1}(S, T) = 0\) (theorem 90).

\textbf{\SGAInlineRefTarget{sga3:xxiv:corollary:4-2-4}Corollary 4.2.4.} \emph{Let \(S\) be a scheme and}

\[ 1 \to G \to G' \to G'' \to 1 \]

\emph{an exact sequence of \(S\)-group schemes (for (fpqc)), \(G\) being semisimple of constant type. Let \(P\) be a principal
homogeneous bundle under \(G'\), suppose the sheaf\footnote{N.D.E.: We have replaced ``bundle'' by ``sheaf'', and then \(G'\) by \(G''\).} \(P / G\) associated representable (for example \(G''\)
affine over \(S\)). For \(P\) to be locally isotrivial (resp. semi-locally isotrivial), it is necessary and sufficient that
\(P / G\) be so (as bundle under \(G''\)).}

If \(P\) is trivial, \(P / G\) is too, which shows that the condition is necessary. Conversely, suppose \(S\) local (resp.
semi-local) and \(P / G\) isotrivial, so trivialized by a finite surjective étale extension \(S'\) of \(S\). Extending the
base to \(S'\), one may reduce the structure group of \(P_{S'}\) to \(G_{S'}\). But \(S'\) is still semi-local and \(G_{S'}\)
semisimple of constant type, so the obtained bundle is isotrivial (\SGARefLink{sga3:xxiv:proposition:4-2-2}{4.2.2}).

\subsection*{\SGAInlineRefTarget{sga3:xxiv:subsection:4-3}4.3. Proof: general case}

Let us first note that \SGARefLink{sga3:xxiv:theorem:4-1-5:item:ii}{4.1.5 (ii)} follows at once from the application of \SGARefLink{sga3:xxiv:corollary:4-2-4}{4.2.4} to the exact sequence

\[ 1 \to d\acute{e}r(G) \to G \to corad(G) \to 1. \]

Let us therefore prove \SGARefLink{sga3:xxiv:proposition:4-2-2:item:i}{(i)}. If \(G\) is split, \(rad(G)\) and \(corad(G)\) are split, as is \(rev(G)\); so \SGARefLink{sga3:xxiv:theorem:4-1-5:item:a}{(a)} implies \SGARefLink{sga3:xxiv:theorem:4-1-5:item:b}{(b)},
(b'), and \SGARefLink{sga3:xxiv:theorem:4-1-5:item:c}{(c)}.

\subsubsection*{4.3.1.}

One has \SGARefLink{sga3:xxiv:theorem:4-1-5:item:c}{(c)} \(\Rightarrow\) \SGARefLink{sga3:xxiv:theorem:4-1-5:item:a}{(a)}. Let \(R\) be the type of \(G\); one has an exact sequence

\[ 1 \to ad(\acute{E}p_{S}(R)) \to A_{S}(R) \to E_{S} \to 1. \]

Applying \SGARefLink{sga3:xxiv:corollary:4-2-4}{4.2.4} to the bundle \(P = Isom_{S-gr.}(\acute{E}p_{S}(R), G)\) and to the associated bundle \(rev(G) = P /
ad(\acute{E}p_{S}(R))\), one has \SGARefLink{sga3:xxiv:theorem:4-1-5:item:c}{(c)} \(\Rightarrow\) \SGARefLink{sga3:xxiv:theorem:4-1-5:item:a}{(a)}.

\subsubsection*{4.3.2.}

One has \SGARefLink{sga3:xxiv:theorem:4-1-5:item:b}{(b)} \(\Rightarrow\) \SGARefLink{sga3:xxiv:theorem:4-1-5:item:a}{(a)}. It suffices to prove that if \(rad(G)\) is split, \(G\) is semi-locally isotrivial. Now let \(G_{0} =
\acute{E}p_{S}(R)\); consider the category of pairs \((G',f)\), where \(G'\) is a form of \(G_0\) and \(f\) is an isomorphism of
\(rad(G_{0})\) onto \(rad(G)\).

One knows (\SGARefLink{sga3:xxiv:2.16}{2.16}) that this category is equivalent to the category of principal homogeneous bundles under a certain
\(S\)-group \(H\) extension of a finite twisted constant group by a semisimple group. It suffices to prove that every
principal bundle under \(H\) is semi-locally isotrivial; this follows at once from \SGARefLink{sga3:xxiv:corollary:4-2-4}{4.2.4}.

\subsubsection*{4.3.3.}

One has (b') \(\Rightarrow\) \SGARefLink{sga3:xxiv:theorem:4-1-5:item:a}{(a)}. One can argue as previously (it will moreover be the same group \(H\) that arises). One can also see
that \SGARefLink{sga3:xxiv:theorem:4-1-5:item:b}{(b)} and (b') are equivalent: a torus isogenous to a locally split torus is also locally split (cf.~\SGARefLink{sga3:ix:proposition:2:11:item:iii:01}{Exp. IX, 2.11
(iii)}).

\subsubsection*{4.3.4.}

One has \SGARefLink{sga3:xxiv:theorem:4-1-5:item:d}{(d)} \(\Rightarrow\) \SGARefLink{sga3:xxiv:theorem:4-1-5:item:a}{(a)}. It suffices to prove that a group of constant type possessing a split maximal torus is semi-locally
isotrivial; this follows at once from \SGARefLink{sga3:xxiv:2.14}{2.14}.

\subsubsection*{4.3.5.}

One has \SGARefLink{sga3:xxiv:theorem:4-1-5:item:a}{(a)} \(\Rightarrow\) \SGARefLink{sga3:xxiv:theorem:4-1-5:item:d}{(d)}. It suffices to prove that a maximal torus of a split group is semi-locally isotrivial. More
precisely:

\textbf{\SGAInlineRefTarget{sga3:xxiv:lemma:4-3-6}Lemma 4.3.6.} \emph{Let \(S\) be a semi-local scheme, \(G\) an \(S\)-reductive group, \(T_0\) a split maximal torus of \(G\), \(W_{0}
= Norm_{G}(T_{0}) / T_{0}\) (this is a locally constant \(S\)-group, constant if \(G\) is of constant type, by \SGARefLink{sga3:xxiv:2.14}{2.14}), \(T\) a
maximal torus of \(G\).}

\emph{There exists a morphism \(S' \to S\) which is principal homogeneous under \(W_0\) (hence étale finite and surjective, and
even principal Galois if \(G\) is of constant type) such that \(T_{S'}\) is conjugate to \((T_{0})_{S'}\) by an element of
\(G(S')\) (and so in particular split).}

Indeed, one knows that \(Transp_{G}(T_{0}, T)\) is a principal homogeneous bundle under \(Norm_{G}(T_{0})\) (cf.~for example
\SGARefLink{sga3:xi:corollary:5:4:bis}{Exp. XI, 5.4 bis}). Set \(S' = Transp_{G}(T_{0}, T) / T_{0}\); this is a principal homogeneous bundle under \(W_0\).
Extending the base from \(S\) to \(S'\), one can reduce the structure group of \(Transp_{G}(T_{0}, T)\) to \(T_0\). Now \(S'\) is
semi-local and \(T_0\) split, so \(Transp_{G}(T_{0}, T)\) possesses a section over \(S'\).\footnote{N.D.E.: by ``theorem 90''.} QED

\subsection*{\SGAInlineRefTarget{sga3:xxiv:subsection:4-4}4.4. Use of the existence of maximal tori}

Using Grothendieck's existence theorem for maximal tori (\SGARefLink{sga3:xiv:3.20}{Exp. XIV, 3.20}), one can considerably sharpen the preceding
results. Let us state at once:

\textbf{\SGAInlineRefTarget{sga3:xxiv:theorem:4-4-1}Theorem 4.4.1.} \emph{Let \(S\) be a semi-local scheme, \(R\) a pinned reduced root datum, \(W\) its Weyl group, \(E\) the group
of its automorphisms. (Recall that \(E\) operates naturally on \(W\) and that the semi-direct product \(A = W \cdot E\) is
identified with the group of automorphisms of \(R\) non-pinned, cf.~\SGARefLink{sga3:XXI:6.7.2}{Exp. XXI, 6.7.2}).}

\emph{(i) Every principal homogeneous bundle under \(\acute{E}p_{S}(R)\) is trivialized by a principal Galois covering \(S' \to
S\) of group \(W\).}

\emph{(ii) Let \(G\) be an \(S\)-reductive group of type \(R\); let \(rev(G) = Isomext(\acute{E}p_{S}(R), G)\) be the associated
Galois covering of \(S\) with group \(E\). Let \(W_0\) be the form of \(W_S\) associated to \(rev(G)\). There exists a morphism
\(S' \to S\), which is a principal homogeneous bundle under \(W_0\), such that \(G_{S'}\) is quasi-splittable (i.e.~possesses
a Borel subgroup).}

\emph{\SGAInlineRefTarget{sga3:xxiv:theorem:4-4-1:item:iii}(iii) Every \(S\)-reductive group of type \(R\) is split by a principal Galois covering \(\tilde{S} \to S\) with group \(A\).}

Let us first state:

\textbf{\SGAInlineRefTarget{sga3:xxiv:proposition:4-4-2}Proposition 4.4.2.} \emph{Let \(S\) be a scheme, \(G\) an \(S\)-reductive group, \(T\) a maximal torus of \(G\), \(P\) a principal
homogeneous bundle under \(G\), \(G'\) the twisted form of \(G\) associated to \(P\) (one then has, cf.~\SGARefLink{sga3:xxiv:proposition:4-2-1}{4.2.1}, a canonical
isomorphism \(P / Norm_{G}(T) \xrightarrow{\sim} Tor(G')\)). Let \(T'\) be a maximal torus of \(G'\), defining a composite
morphism}

\[
  S\longrightarrow\operatorname{Tor}(G')
  \xrightarrow{\sim}P/\operatorname{Norm}_G(T).
\]

\emph{Consider the canonical morphisms}

\[
  P\longrightarrow P/T\longrightarrow P/\operatorname{Norm}_G(T).
\]

\emph{and take their inverse images by the preceding morphism:}

\[
  \begin{tikzcd}[column sep=large,row sep=large]
P \arrow[r]
&
P/T \arrow[r]
&
P/\operatorname{Norm}_G(T)
\\
H \arrow[u] \arrow[r]
&
S' \arrow[u] \arrow[r]
&
S \arrow[u]
\end{tikzcd}

\]

\emph{Then \(S'\) (resp. \(H\)) is a principal homogeneous bundle over \(S\) (resp. \(S'\)) under \(W_{G}(T)\) (resp. \(T_{S'}\)).
Moreover, if one makes \(W_{G}(T)\) operate on \(T\) in the obvious way, the bundle associated to \(S'\) is isomorphic to
\(T'\).}

The first two assertions are trivial, the last is proved like the corresponding assertion of \SGARefLink{sga3:xxiv:2.6}{2.6}.

\textbf{\SGAInlineRefTarget{sga3:xxiv:corollary:4-4-3}Corollary 4.4.3.} \emph{Let \(S\) be a semi-local scheme, \(G\) an \(S\)-reductive group, \(T\) a maximal torus of \(G\). Suppose
one of the two following conditions is satisfied:}

\emph{\SGAInlineRefTarget{sga3:xxiv:corollary:4-4-3:item:i}(i) \(T\) is split.}

\emph{\SGAInlineRefTarget{sga3:xxiv:corollary:4-4-3:item:ii}(ii) \(T\) is contained in a Borel subgroup of \(G\), and \(G\) is either adjoint, or simply connected.}

\emph{Let moreover \(P\) be a principal homogeneous bundle under \(G\). There exists an \(S\)-scheme \(S'\), which is a principal
homogeneous bundle under \(W_{G}(T)\), such that \(P_{S'}\) is trivial.}

Indeed, if \(G'\) is the form of \(G\) associated to \(P\), then \(G'\) possesses a maximal torus \(T'\) (\SGARefLink{sga3:xiv:3.20}{Exp. XIV, 3.20}).
Resuming the notation of the preceding proposition, one sees that \(H^{1}(S', T_{S'}) = 0\) (by theorem 90 for \SGARefLink{sga3:xxiv:corollary:4-4-3:item:i}{(i)}, by
\SGARefLink{sga3:xxiv:corollary:3-14}{3.14} for \SGARefLink{sga3:xxiv:corollary:4-4-3:item:ii}{(ii)}). The morphism \(H \to S'\) possesses a section, so \(P_{S'}\) also possesses a section over \(S'\). QED

Let us now prove the theorem. Assertion \SGARefLink{sga3:xxiv:corollary:4-4-3:item:i}{(i)} is a particular case of the preceding corollary (take \(G =
\acute{E}p_{S}(R)\), endowed with its canonical split torus). Let us prove \SGARefLink{sga3:xxiv:corollary:4-4-3:item:ii}{(ii)}. One knows (\SGARefLink{sga3:xxiv:3.12}{3.12}) that \(G\) is an inner
twisted form of

\[ G_{0} = Q-\acute{E}p_{rev(G)/S}(R). \]

If \(T_0\) is the canonical maximal torus of \(G_0\), \(W_{G_{0}}(T_{0})\) is the group \(W_0\) described in the statement. The
form \(G\) of \(G_0\) corresponds to a principal homogeneous bundle \(P\) under \(ad(G_{0})\) (\(P = Isomint(G_{0}, G)\)). The
group \(W_{ad(G_{0})}(T^{ad}_{0})\) is canonically isomorphic to \(W_0\), and one obtains the wanted result by applying
\SGARefLink{sga3:xxiv:corollary:4-4-3}{4.4.3} to the situation \((ad(G_{0}), T^{ad}_{0}, P)\), hypothesis \SGARefLink{sga3:xxiv:corollary:4-4-3:item:ii}{(ii)} of \SGARefLink{sga3:xxiv:corollary:4-4-3}{4.4.3} being well verified. Let us finally prove
\SGARefLink{sga3:xxiv:theorem:4-4-1:item:iii}{(iii)}. Resume the notation of \SGARefLink{sga3:xxiv:corollary:4-4-3:item:ii}{(ii)}; one has a diagram

\[
  \begin{tikzcd}[column sep=large,row sep=large]
&
\operatorname{rev}(G) \arrow[d,"E_S"']
\\
S'
  \arrow[r,"W_0"']
&
S
\end{tikzcd}

\]

One knows that \(G_{S'}\) is isomorphic to \((G_{0})_{S'}\) and that \((G_{0})_{rev(G)}\) is splittable. If one sets
\(\tilde{S} = S' \times_{S} rev(G)\), \(G_{\tilde{S}}\) is splittable, and it only remains to verify that \(\tilde{S}\) is
indeed a principal Galois covering of \(S\) with group \(A\), which follows from the following more general lemma (naturally
valid in any site):
